\scanpage{117}
% ms. line 01
Montrons qu'alors $\frac{n(r)}{r^\alpha}\to0$ ; si en effet on avait
% ms. line 02
$\frac{n(b_i)}{b_i^\alpha}>\varepsilon$ pour des $b_i\to+\infty$, on aurait
% ms. line 03
$n(b_i)-n(b_{i-1})\geq\varepsilon\lambda_i^\alpha-M\lambda_{i-1}^\alpha$ d'où
% ms. line 04
\[ \sum_{\substack{\unsure{\nu|v}\\b_{i-1}<r_\nu\leq b_i}}r_\nu^{-\alpha}\geq\big(n(b_i)-n(b_{i-1})\big)\lambda_i^{-\alpha}\geq\varepsilon-M\left(\frac{\lambda_{i-1}}{\lambda_i}\right)^\alpha \]
% ms. line 05
on peut supposer $M\left(\frac{\unsure{b_{i-1}|\lambda_{i-1}}}{\lambda_i}\right)^\alpha\leq\frac{\varepsilon}{2}$ pour tout $i$,
% ms. line 06
d'où $\sum_{b_{i-1}<r_\nu\leq b_i}r_\nu^{-\alpha}\geq\frac{\varepsilon}{2}$, \deleted{d'où $\sum$} ce
% ms. line 07
qui est contradictoire.

% ms. line 08
Si alors on fait $b\to\infty$ dans (1 bis), on
% ms. line 09
obtient (2).

% ms. line 10
\deleted{\illegible{1}} Des résultats analogues seraient obtenus si
% ms. line 11
on considérait les sommes $\sum_\nu\lambda(r_\nu)$, où
% ms. line 12
$\lambda$ est fonction définie sur $R^{+*}$, continûment
% ms. line 13
différentiable, \inserted{décroissante,} et tendant vers zéro
% ms. line 14
\deleted{\unsure{pour}} à l'infini\unsure{\}} : Alors
% ms. line 15
\[ \boxed{\displaystyle \sum_\nu\lambda(r_\nu)<+\infty\ \Longleftrightarrow\left\{\begin{aligned}&\int_0^\infty n(r)[-\lambda'(r)]\,dr<+\infty\\[-2pt]&\deleted{\text{et }\lambda(r)n(r)\ \text{\unsure{reste borné}}}\end{aligned}\right.} \]
% ms. line 16
et alors $\lambda(r)n(r)\to0$ pour $r$ infini, et
% ms. line 17
on a \quad $\displaystyle\sum_{r_\nu>a}\lambda(r_\nu)=-\lambda(\deleted{\unsure{b}}a)n(a+0)+\int_a^\infty n(r)(-\lambda'(r))\,dr\deleted{\unsure{]}}$
% ms. line 18
\deleted{\unsure{Ceci}}
% archive sheet number: 118
